\section{El recorrido, momento por momento}
El trámite avanza en cuatro momentos: 1. Reúne tus requisitos, 2. Propón sínodo y fecha, 3. Invitaciones al sínodo, 4. Examen y cierre. Es igual en todas las modalidades; las pantallas son de la modalidad \emph{Tesis}.
\subsection{Dar el visto bueno}
Contesta las tres preguntas (tesis concluida, entregada a usted, revisión en el periodo), escribe la vigencia de tu firma electrónica UAEMéx, marca que leíste el documento y firma.
\captura{tesis_sinodal_01_visto_bueno}{Visto bueno de un integrante del sínodo desde su enlace.}
\captura{tesis_sinodal_02_visto_bueno_externo}{Visto bueno del sinodal externo (sin cuenta): contesta tres preguntas, declara la vigencia de su firma electrónica UAEMéx y firma.}
\captura{tesis_sinodal_06_visto_bueno_cuenta}{El visto bueno con sesión: el mismo formulario, con la vigencia de su perfil ya escrita.}
\subsection{Si tu firma UAEMéx no alcanza}
El enlace te avisa y te deja actualizar la fecha; el egresado ve el aviso en su siguiente paso.
\captura{tesis_sinodal_03_vigencia_no_alcanza}{Si la vigencia declarada no cubre un mes después del examen, el enlace lo avisa y permite actualizar la fecha.}
\captura{tesis_sinodal_04_vigencia_actualizar}{El sinodal actualiza la vigencia desde el mismo enlace.}
\subsection{Con cuenta: «Posgrado: firmas»}
Si tienes cuenta con el mismo correo, tus firmas pendientes y hechas están en \emph{Posgrado: firmas}.
\captura{tesis_sinodal_05_mis_firmas}{Con cuenta en Traza FI, el sinodal encuentra sus firmas en «Posgrado: firmas».}
\subsection{Acusar de recibido tu invitación}
Con la invitación firmada por la Coordinación recibes un enlace: consulta tu oficio y pulsa \boton{Acusar de recibido}.
\captura{tesis_sinodal_07_acuse}{Invitación al sínodo: el integrante consulta su oficio y acusa de recibido en línea.}
\subsection{Después del examen}
Tus firmas hechas quedan en \emph{Posgrado: firmas}; tus constancias las emite Apoyo a Grado y las firma la Coordinación.
\captura{tesis_sinodal_08_firmadas}{En «Posgrado: firmas» quedan las firmas hechas, para consultar el documento con tu sello.}
\subsection{La vigencia en «Mi perfil»}
Con cuenta, declara la vigencia de tu firma electrónica UAEMéx en \emph{Mi perfil}; la plataforma te avisa cuando le falte un mes o menos.
\captura{tesis_sinodal_09_perfil}{«Mi perfil»: la vigencia de tu firma electrónica UAEMéx y el aviso cuando le falta un mes o menos.}
\captura{tesis_sinodal_10_aviso_vencimiento}{El aviso de vencimiento también llega a tus notificaciones.}
